\begin{helper}[FiberAndDegreeSameColorLiftedIntersectionSharpBaseNegligible]
\label{helper:FiberAndDegreeSameColorLiftedIntersectionSharpBaseNegligible}
Let \(\beta=1/2\), let
\[
m(n)=\noderef{TwoBiteNaturalReducedVertexCount}(n),
\qquad
p(n)=\noderef{TwoBiteEdgeProbability}(\beta,n),
\]
and take probabilities with respect to
\(\noderef{TwoBiteEventProbability}(n,m(n),p(n))\).  The event that a
same-color lifted-neighborhood intersection is larger than
\(100(\log n)^3\) while every same-color base codegree is at most \(\log n\)
has \(\noderef{NegligibleProbability}\).  That is, the negligible event is
\[
\begin{aligned}
&\Bigl[
\exists r\ne s,\,
|\noderef{TwoBiteLiftedNeighborhood}(\omega,\operatorname{inl}(r))
\cap
\noderef{TwoBiteLiftedNeighborhood}(\omega,\operatorname{inl}(s))|
>100(\log n)^3\\
&\quad\text{or}\quad
\exists b\ne c,\,
|\noderef{TwoBiteLiftedNeighborhood}(\omega,\operatorname{inr}(b))
\cap
\noderef{TwoBiteLiftedNeighborhood}(\omega,\operatorname{inr}(c))|
>100(\log n)^3
\Bigr]\\
&\quad\cap
\Bigl[
\forall r\ne s,\,
|\{t:\noderef{TwoBiteRedGraph}(\omega).Adj(r,t)\land
\noderef{TwoBiteRedGraph}(\omega).Adj(s,t)\}|\le\log n,\\
&\qquad\quad
\forall b\ne c,\,
|\{t:\noderef{TwoBiteBlueGraph}(\omega).Adj(b,t)\land
\noderef{TwoBiteBlueGraph}(\omega).Adj(c,t)\}|\le\log n
\Bigr].
\end{aligned}
\]
\end{helper}

\begin{proof}
Write \(M=m(n)\), \(p=p(n)\), and \(L=\log n\).  We discard finitely many
\(n\) so that \(L\ge1\), \(M=\lfloor n/L^2\rfloor\), \(n\le M^2\),
\(n\le 2L^2M\), and \(0\le p\le1\).  These eventual parameter facts follow
from \(\noderef{TwoBiteNaturalReducedVertexCount}\),
\(\noderef{TwoBiteReducedVertexCount}\), \(\noderef{TwoBiteStretch}\),
\(\noderef{TwoBiteNaturalTotalMassOneEventually}\), and
\(\noderef{OppositeProjectionEdgeProbBounds}\).

Fix base graphs \(G_R,G_B\) on \(\operatorname{Fin}(M)\).  A sample with
these base graphs is determined by its injection
\[
\pi:\operatorname{Fin}(n)\hookrightarrow
\operatorname{Fin}(M)\times\operatorname{Fin}(M)
\]
from \(\noderef{TwoBiteSample}\).  Expanding
\(\noderef{TwoBiteEventProbability}\) and
\(\noderef{TwoBiteSampleWeight}\), the contribution of an injection event
\(F\) inside this base-graph cell is
\[
\sum_{\pi\in F}
\noderef{GnpGraphWeight}(p,G_R)\,
\noderef{GnpGraphWeight}(p,G_B)\,
\noderef{UniformInjectionWeight}(n,M).
\tag{1}
\]
The denominator for the whole base-graph cell is the same expression with
\(F\) equal to all injections.  If this denominator is \(0\), then the cell
has zero mass and contributes nothing to the unconditional probability.  If it
is nonzero, the two graph-weight factors are constant and cancel in the
quotient.  Since \(n\le M^2\), \(\noderef{UniformInjectionWeight}(n,M)\) is
the reciprocal of the number of injections, so the conditional probability
inside the cell is exactly
\[
\frac{|F|}
{|\{\pi:\operatorname{Fin}(n)\hookrightarrow
\operatorname{Fin}(M)\times\operatorname{Fin}(M)\}|}.
\tag{2}
\]

Assume now that the fixed base graphs satisfy the sharp same-color base
codegree event.  Suppose first that \(r,s\) are distinct red vertices and let
\[
A_{r,s}=\{t:\noderef{TwoBiteRedGraph}(\omega).Adj(r,t)\land
\noderef{TwoBiteRedGraph}(\omega).Adj(s,t)\}
\]
be their red base common-neighbor set.  In the fixed cell this is the same set
computed in \(G_R\), and the sharp event gives \(|A_{r,s}|\le L\).  By the
definitions of
\(\noderef{TwoBiteLiftedNeighborhood}\), \(\noderef{TwoBiteRedProjection}\),
and \(\noderef{TwoBiteEmbedding}\), a lifted vertex belongs to both lifted
red neighborhoods exactly when its injected product cell lies in
\[
A_{r,s}\times \operatorname{Fin}(M).
\]
Thus the number of marked product cells is at most \(LM\).  The hypotheses
of \(\noderef{FiberAndDegreeSameColorLiftedIntersectionFixedPairInjectionTail}\)
therefore apply to this fixed pair.  Combined with the uniform-injection
identity (2), the conditional probability that this red pair has lifted
intersection exceeding \(100L^3\) is at most \(\exp(-2L^3)\).

For distinct blue vertices \(b,c\), the marked cells are
\[
\operatorname{Fin}(M)\times
\{t:\noderef{TwoBiteBlueGraph}(\omega).Adj(b,t)\land
\noderef{TwoBiteBlueGraph}(\omega).Adj(c,t)\},
\]
using \(\noderef{TwoBiteBlueProjection}\), and the sharp base-codegree event
again gives at most \(LM\) marked cells.  This set is in the second coordinate
rather than the first, but the coordinate-swap bijection
\[
\sigma(x,y)=(y,x)
\]
of \(\operatorname{Fin}(M)\times\operatorname{Fin}(M)\) sends injections whose
images meet \(\operatorname{Fin}(M)\times A_{b,c}\) in more than \(100L^3\)
points bijectively to injections whose swapped images meet
\(A_{b,c}\times\operatorname{Fin}(M)\) in more than \(100L^3\) points.  Hence
the count of bad blue injections is the same as the first-coordinate count
controlled by
\(\noderef{FiberAndDegreeSameColorLiftedIntersectionFixedPairInjectionTail}\).
The same conditional bound \(\exp(-2L^3)\) follows for each fixed blue pair.

There are fewer than \(M^2\) ordered red pairs and fewer than \(M^2\) ordered
blue pairs.  Summing the fixed-pair conditional bounds gives, for every
base-graph cell satisfying the sharp event, the conditional bound
\[
2M^2\exp(-2L^3)
\]
for the lifted-intersection bad event.  If the fixed base graphs do not
satisfy the sharp event, then the event in the statement is false in that
cell, because the statement includes the sharp base-codegree assumptions.

It remains to remove the conditioning.  By
\(\noderef{OppositeProjectionEdgeProbBounds}\) and
\(\noderef{TwoBiteSampleWeightNonnegative}\), all two-bite sample weights are
nonnegative eventually.  Summing the cellwise bound over all \(G_R,G_B\) gives
\[
\Pr(\text{event in the statement})
\le
2M^2\exp(-2L^3)
\sum_{G_R,G_B}
\noderef{GnpGraphWeight}(p,G_R)\noderef{GnpGraphWeight}(p,G_B).
\tag{3}
\]
The double graph sum factors into the product of the red and blue graph sums,
and each factor is \(1\) by \(\noderef{GnpGraphWeightSumOne}\).  Therefore
\[
\Pr(\text{event in the statement})\le 2M^2\exp(-2L^3).
\tag{4}
\]

Finally, \(\noderef{OppositeProjectionAsymptoticBound}\), with \(c=2\) and
\(m(n)=\noderef{TwoBiteNaturalReducedVertexCount}(n)\), gives
\[
M(n)^2\exp(-2(\log n)^3)\to0.
\]
The preceding probability bound is a constant multiple of this sequence, so
the event has negligible probability.
\end{proof}
